\documentclass[../../main.tex]{subfiles}
\graphicspath{{../../image/}{../../figure/}}
\begin{document}

\section*{习题}\label{section:ch13-习题}
\addcontentsline{toc}{section}{习题}

\textbf{13.1}（$\star$）证明：对应于\cref{fig:ch13-adjacency}中两种节点排序选择的置换 $(A,B,C,D,E)\rightarrow(C,E,A,D,B)$，可以表示为式\eqref{align:ch13-a-perm}定义的形式，其中置换矩阵由式\eqref{align:ch13-perm-matrix}给出。

\textbf{13.2}（$\star\star$）证明：对应于\cref{fig:ch13-adjacency}中两种节点排序选择的图 $\mathcal{G}$ 中与每个节点相连的边数可以由矩阵 $\bm{A}^2$ 中相应的对角元素给出，其中 $\bm{A}$ 是邻接矩阵。

\textbf{13.3}（$\star$）由下式给出的邻接矩阵，画出其对应的图：
\begin{align}
\bm{A}=\begin{pmatrix}0&1&1&0&1\\1&0&1&1&1\\1&1&0&1&0\\0&1&1&0&0\\1&1&0&0&0\end{pmatrix}.
\label{align:ch13-ex-a}
\end{align}

\textbf{13.4}（$\star\star$）证明在数据矩阵 $\bm{X}$ 前乘以式\eqref{align:ch13-perm-general}定义的置换矩阵 $\bm{P}$ 的效果是创建一个新的数据矩阵 $\widetilde{\bm{X}}$[由式\eqref{align:ch13-x-perm}给出]，其行已根据置换函数 $\pi(\cdot)$ 进行置换。

\textbf{13.5}（$\star\star$）证明转换后的邻接矩阵 $\widetilde{\bm{A}}$，它由式\eqref{align:ch13-a-perm}定义，其中 $\bm{P}$ 由式\eqref{align:ch13-perm-general}定义，使得新矩阵的行和列都根据置换函数 $\pi(\cdot)$ 相对于原始邻接矩阵 $\bm{A}$ 进行置换。

\textbf{13.6}（$\star\star$）将式\eqref{align:ch13-update-mlp}所示的更新方程写成使用矩阵的图级别方程。为了保持符号简洁，省略层索引 $l$。为此，首先将节点嵌入向量 $\bm{h}_n^{(l)}$ 收集到一个矩阵 $\bm{H}$ 中，其行 $n$ 由 $\bm{h}_n^{\mathsf T}$ 给出。然后证明下式
\begin{align}
\bm{Z}=\bm{A}\bm{H}
\label{align:ch13-ex-z}
\end{align}
给出的邻域聚合向量 $\bm{z}_n$ 可以以矩阵形式写成 $\bm{Z}=\bm{A}\bm{H}$，$\bm{Z}$ 是一个 $N\times D$ 大小的矩阵，其中的行 $n$ 由 $\bm{z}_n^{\mathsf T}$ 给出。最后证明式\eqref{align:ch13-update-mlp}中非线性激活函数的参数可以以矩阵形式写成
\begin{align}
\bm{H}^{(l+1)}=f\!\left(\bm{H}\bm{W}_{\mathrm{self}}^{\mathsf T}+\bm{A}\bm{H}\bm{W}_{\mathrm{neigh}}^{\mathsf T}+\bm{1}\bm{b}^{\mathsf T}\right)
\label{align:ch13-ex-mat}
\end{align}
其中 $\bm{1}$ 是一个 $D$ 维的列向量，其中的所有元素都是 $1$。

\textbf{13.7}（$\star\star$）利用深度图卷积网络第 $l$ 层的等变性[式\eqref{align:ch13-layer-equiv}]和节点变量的置换性质[式\eqref{align:ch13-x-perm}]，证明式\eqref{align:ch13-layers}定义的完整深度图卷积网络也是等变的。

\textbf{13.8}（$\star\star$）证明在式\eqref{align:ch13-attn-agg}定义的聚合函数中[注意力权重由式\eqref{align:ch13-attn-mlp}给出]，该聚合函数在图中节点重新排序下是等变的。

\textbf{13.9}（$\star$）证明一个图注意力网络（其中图是全连接的，即每对节点之间都有一条边），等同于标准的~Transformer~架构。

\textbf{13.10}（$\star\star$）当坐标系统被平移时，对坐标系统定义的物体的位置向量使用
\begin{align}
\bm{r}'=\bm{r}+\bm{c}
\label{align:ch13-ex-trans}
\end{align}
进行转换，其中 $\bm{c}$ 是描述平移的固定向量。类似地，如果坐标系统发生旋转和/或镜像反射，则对坐标系统定义的物体的位置向量使用
\begin{align}
\bm{r}'=\bm{R}\bm{r}
\label{align:ch13-ex-rot}
\end{align}
进行转换，其中 $\bm{R}$ 是一个正交矩阵，它的逆由其转置给出，所以
\begin{align}
\bm{R}\bm{R}^{\mathsf T}=\bm{R}^{\mathsf T}\bm{R}=\bm{I}.
\label{align:ch13-ex-orth}
\end{align}
利用这些属性，证明在平移、旋转和反射下，消息传递方程[式\eqref{align:ch13-geo-edge}、式\eqref{align:ch13-geo-agg}和式\eqref{align:ch13-geo-node}]是不变的，并且式\eqref{align:ch13-geo-coord}给出的坐标嵌入向量是等变的。

\end{document}
